\section{Introduction}
\label{sec:intro}

Language models learn rich distributional structure during pretraining---including
the ability to produce text spanning many registers, genres, and authorial voices.
Alignment procedures such as RLHF~\cite{ouyang2022instructgpt} and
DPO~\cite{rafailov2023dpo} reshape these distributions toward human preferences,
yielding models that are helpful, harmless, and honest. But a persistent concern
accompanies this process: \emph{does alignment narrow the model's expressive range?}
\citet{kirk2023rlhf} showed that RLHF causes measurable mode collapse in both
per-input and across-input diversity. \citet{lin2023mitigating} documented
degradation on downstream benchmarks, coining the term ``alignment tax.'' If
alignment suppresses stylistic distributions that existed in the base model, then
creative applications---authorial impersonation, genre writing, code-switching
across registers---may suffer.

{\bf But is this the right framing?} The Superficial Alignment
Hypothesis~\cite{zhou2023lima} suggests that alignment is a thin veneer atop
pretraining distributions, adjusting format and tone without fundamentally altering
the model's knowledge or capabilities. \citet{ji2024resistance} provided further
evidence: aligned models exhibit ``elasticity,'' reverting toward pretraining
distributions upon further fine-tuning. If alignment is superficial and reversible,
then the base model's stylistic distributions should survive alignment intact---and
distribution interpolation should be able to navigate continuously between base and
aligned behaviors.

{\bf Our approach.} We directly test these ideas by measuring stylistic
impersonation across seven distinct literary and functional styles (Hemingway,
Shakespeare, Poe, Austen, Tolkien, academic, noir) using paired base and aligned
models (\qwenbase and \qweninst). We apply logit-space distribution arithmetic---
linearly interpolating base and aligned logits at each decoding step---at five
interpolation coefficients ($\interpcoeff \in \{0, 0.25, 0.5, 0.75, 1\}$). A
\gptjudge judge evaluates style fidelity and text quality, supplemented by
automated diversity metrics (Distinct-N, Self-BLEU).

{\bf What we find.} The results overturn the naive alignment tax hypothesis.
The aligned model \emph{significantly outperforms} the base model on
style fidelity ($7.26$ vs.\ $3.74$, $p < 0.0001$) and is the only
condition that reliably differentiates between target styles (ANOVA
$F = 6.32$, $p = 0.002$; base: $F = 0.87$, $p = 0.54$). Distribution
arithmetic produces smooth, monotonic interpolation, confirming that base
distributions survive alignment. But mixing in base logits \emph{reduces}
style quality rather than recovering it, because the base model lacks the
instruction-following mechanism needed to address its own distributions.

In summary, our main contributions are:
\begin{itemize}[leftmargin=*,itemsep=0pt,topsep=0pt]
    \item We conduct the first controlled comparison of stylistic impersonation
    ability before and after alignment, finding that alignment \emph{improves}
    rather than diminishes controllable access to diverse styles.
    \item We apply logit-space distribution arithmetic to style generation,
    demonstrating smooth interpolation between base and aligned behavior across
    seven styles and three topics (231 total generations).
    \item We reframe the alignment--diversity tradeoff: the relevant axis is not
    preservation versus loss of distributions, but controllability versus
    distributional freedom. Alignment adds an addressing mechanism that unlocks
    latent stylistic diversity.
\end{itemize}
